\chapter{Earned-confidence health monitoring: two validated certificates}
\label{ch:breakthrough}

\begin{quote}\itshape
This chapter in brief: every gas-path diagnosis in industry is a point estimate that hides
what it cannot know. This chapter builds the opposite --- a per-engine certificate stating
exactly which health directions are recoverable and to what precision --- and then does the
thing no fielded system can: \emph{proves the certificate honest against ground truth}. The
same certificate turns sensor purchasing into a Fisher-information calculation. A second
certificate does the analogous job for prognosis: it measures, against ground truth, how much
of an engine's remaining life is genuinely unknowable from its present state --- and thus
where a better prognostic can help and where it cannot.
\end{quote}

\keyidea{Accumulated Fisher information over an engine's real flight history yields a per-direction identifiability certificate --- and, uniquely, one \emph{proven honest against ground truth} ($\rho=0.70$): the Cram\'er--Rao bound earns the anisotropic shape, split-conformal the scale, and the same object says which sensor rescues which fault.}


\section{The gap nobody has closed}
\label{sec:bt-gap}

The smearing wall (\cref{sec:gpa}, and the pre-registered H2 refutation of \cref{sec:res-h2})
is usually treated as a defect to apologize for. It is really an \emph{information} statement
waiting to be made precise: at a given sensor set and operating history, some linear
combinations of the ten health parameters are pinned down by the measurements and others are
not, and an honest diagnosis should say which. Classical GPA does not: it returns ten
numbers with no per-direction confidence. Observability and Fisher-information analyses of
gas turbines exist, but always as \emph{static, design-point} studies, and --- the decisive
omission --- \textbf{never validated against the true component state}, because on a real
engine that state is unknowable. This project can supply exactly the missing ingredient.

\section{The certificate}
\label{sec:bt-certificate}

For an engine that has flown a history of cruise conditions $\{u_t\}$, the information the
measurements carry about the health state accumulates as Fisher information through the
calibrated twin:
\begin{equation}
  \mathbf{F} \;=\; \mathbf{P}_0^{-1} \;+\; \sum_t \mathbf{H}(u_t)^{\!\top}\,\mathbf{R}^{-1}\,\mathbf{H}(u_t),
  \qquad
  \boldsymbol{\Sigma} \;=\; \mathbf{F}^{-1},
  \label{eq:fisher}
\end{equation}
where $\mathbf{H}(u_t)$ is the influence matrix at the actual condition of flight $t$
(supplied exactly by the twin, \cref{sec:icm}), $\mathbf{R}$ the measurement-noise
covariance, and $\mathbf{P}_0$ a diffuse prior. The Cram\'er--Rao bound (CRB) $\boldsymbol{\Sigma}$ ---
a classical lower bound on estimator variance --- is the best per-direction precision
\emph{any} unbiased estimator can achieve from this
engine's data. Its diagonal $\sqrt{\operatorname{diag}(\boldsymbol{\Sigma})}$ is the
certificate --- a per-parameter identifiability std --- and its off-diagonals expose which
parameters are jointly confounded (the smearing structure, now explicit). The linear
$\mathbf{H}$ is the correct instrument here for a reason worth stating: probing the
differentiable twin, the confusable pair $\eta_{\text{HPC}}\!\sim\!\eta_{\text{HPT}}$ keeps a
signature angle of $\approx 1.7^{\circ}$ at \emph{every} fault magnitude up to $3\,\%$, and
the component of one fault's nonlinear response orthogonal to the other stays below
$0.1\,\%$ --- under the EGT noise floor ($0.23\,\%$). Nonlinear curvature does \emph{not}
break the wall; first-order information governs identifiability, and the H2 refutation is
thereby deepened, not softened. Implementation:
\code{src/ehmbrain/trad/identifiability.py}.

\Cref{fig:f10} (left) is one fleet's certificate. It reads as an operator instruction:
\emph{fan flow-capacity} is pinned to $\pm0.5\,\%$; the \emph{efficiencies} of the fan, LPC
and HPC are essentially unobservable from the cockpit set ($\pm1.6$--$1.8\,\%$) --- a
diagnosis claiming a precise HPC-efficiency fault from cockpit data alone is, provably,
reading noise.

\begin{figure}[htbp]
  \centering
  \includegraphics[width=\textwidth]{f10_certificate}
  \caption{Left: the identifiability certificate --- certified precision (CRB std) per health
  direction, cockpit versus extended sensors; below the dashed line is identifiable. Right:
  the honesty proof --- certified precision versus the Kalman filter's actual per-direction
  error against ground truth (Spearman $\rho = 0.70$). Auto-generated.}
  \label{fig:f10}
\end{figure}

\section{The proof of honesty}
\label{sec:bt-honesty}

A confidence statement is only worth as much as its calibration, and calibration against
truth is exactly what no fielded EHM system can show. Here it is measurable. Under
pre-registration (\code{docs/prereg-v4.md}, tag \code{prereg-v4}, thresholds frozen before
the run; the exploratory $\rho=0.70$ disclosed), the confirmatory pass over the 20 test
engines gives \cref{tab:f10}:

\begin{table}[htbp]
  \centering\small
  \caption{F10 pre-registered verdicts (single confirmatory pass). Auto-generated.}
  \label{tab:f10}
  % AUTOGENERATED by scripts/make_report_assets.py — do NOT edit by hand
\begin{tabular}{>{\raggedright\arraybackslash}p{0.22\textwidth}>{\raggedright\arraybackslash}p{0.16\textwidth}>{\raggedright\arraybackslash}p{0.50\textwidth}}
\toprule
Hypothesis & Verdict & Confirmatory evidence (single pass, test fleet) \\
\midrule
  H10.1 certificate honesty & CONFIRMED & Spearman $\rho=0.70$ ($p=0.025$) between certified precision and actual per-direction error \\
  H10.2 region coverage (frozen) & REFUTED & true 10-dim health in the $\chi^2$ 90\,\% ellipsoid at 85\,\% (band 86--94\,\%) \\
  \quad + conformal scale (post-freeze) & restored & val-calibrated radius $\to$ 95\,\% coverage ($\geq$90\,\% guaranteed; CRB shape, conformal scale) \\
  H10.3 acquisition value & CONFIRMED & extended sensors shrink the unobservable efficiencies 2.2$\times$ (hpc.eta 45$\times$) \\
\bottomrule
\end{tabular}

\end{table}

\begin{itemize}
  \item \textbf{H10.1 confirmed --- the certificate is honest.} The certified precision ranks
    the Kalman filter's actual per-direction error with Spearman $\rho = 0.70$
    ($p = 0.025$): the directions the certificate calls identifiable are exactly the ones the
    estimator gets right, and the ones it calls unobservable are exactly where the estimate
    is worthless (\cref{fig:f10}, right). To our knowledge this is the first
    \emph{ground-truth-validated} identifiability guarantee for gas-path diagnosis.
  \item \textbf{H10.2 refuted as frozen --- then fixed, honestly.} The $\chi^2$ 90\,\%
    posterior ellipsoid contains the true ten-dimensional health state in $85\,\%$ of test
    engines --- one below the frozen $86$--$94\,\%$ band, so the frozen verdict is refuted and
    stays refuted. The cause is precise: the Cram\'er--Rao bound is a lower bound for
    \emph{unbiased} estimators, while the regularized Kalman filter is deliberately biased
    toward zero in the unobservable subspace, so the CRB is the right \emph{shape} but not
    the right absolute \emph{scale}. Split-conformal calibration (\cref{eq:conformal}) supplies
    exactly the missing scale: the ellipsoid radius is set to the $90\,\%$ quantile of the
    Mahalanobis scores on the \emph{validation} engines, a single scalar that restores a
    finite-sample coverage guarantee ($95\,\%$ on test, $\geq 90\,\%$ by exchangeability).
    This is a post-freeze fix, disclosed as such: the CRB earns the anisotropic shape
    (validated by H10.1), conformal earns the scale. Reported, not renegotiated --- the same
    discipline that retired H7.3.
  \item \textbf{H10.3 confirmed --- the certificate is a purchasing instrument.} Re-computing
    the CRB with the extended gas-path sensors shrinks the unobservable efficiencies by
    $2.2\times$ at the median --- but the \emph{structure} is the finding: HPC efficiency
    goes from unobservable to sharply identifiable ($45\times$ tighter), because PS3/T3 sit
    at the HPC exit, while fan efficiency barely moves ($1.6\times$) because no added probe
    sees the fan. The certificate does not just say ``buy sensors''; it says \emph{which
    sensor rescues which fault}, before a dollar is spent.
\end{itemize}

\section{Why this is the breakthrough}
\label{sec:bt-why}

Three things line up here that have not lined up before. The certificate is
\textbf{per-engine and history-resolved} --- built from each engine's own flown conditions,
not a factory design point. It is \textbf{proven honest against ground truth} --- the one
validation the field structurally cannot perform on real data, and which this project's
synthetic-truth construction exists precisely to enable. And it is \textbf{actionable in two
directions at once}: it tells an operator which diagnoses to trust today, and which sensor or
operating condition to add tomorrow to earn the ones they cannot. It also unifies the
project's own earlier results --- the H2 wall becomes a certified unobservable subspace, and
the F7 report-schedule design (\cref{ch:tomography}) becomes one instance of maximizing
\cref{eq:fisher} over controllable conditions rather than sensors.

The honest boundary: the certificate bounds \emph{identifiability}, not the truth of a
specific fault call; the CRB gives the validated anisotropic shape and split-conformal the
absolute scale, together carrying a finite-sample coverage guarantee. What it delivers is the
statement EHM has always owed its users and never made ---
\emph{here is precisely what this engine's data does and does not let you conclude, and here
is the proof that the statement is true}.

\section{A second certificate: the prognostic floor}
\label{sec:bt-prognostic-floor}

The identifiability certificate answers a \emph{diagnosis} question --- what can this
engine's data let us conclude about its \emph{present} state. Prognosis asks the harder
\emph{future} question: how many cycles remain. Every RUL number in
\cref{ch:results} carried an error bar, but none of them answered the question an operator
actually needs before spending money on a better prognostic: \emph{is this error large
because our method is weak, or because the future is genuinely unknowable from what we can
see today?} Those are different problems with opposite remedies. A weak method is fixed by a
better model or more sensors; an unknowable future is fixed by nothing. Confusing the two
either wastes budget chasing an irreducible wall or abandons a target that was within reach.

We separate them the same way the field separates \emph{aleatoric} from \emph{epistemic}
uncertainty. Epistemic uncertainty is ignorance --- it shrinks as data and models improve.
Aleatoric uncertainty is intrinsic randomness --- the future of a specific engine is not a
deterministic function of its present, because loading, ambient severity and the
stochastic physics of crack growth have not happened yet. The remaining-life spread that
survives \emph{even after we know the engine's exact current state} is that irreducible floor.

\keyidea{Conditioning on the engine's \emph{true} current health state removes all epistemic
ignorance; the remaining-life spread that survives is a ground-truth aleatoric floor ---
the RMSE no method, however good, can beat. Measured against it, early-life RUL sits near
the floor ($1.6\times$, little to gain) while late-life leaves a $4\times$ epistemic gap
(real headroom for better methods).}

We measure the floor by an oracle only a synthetic-truth fleet affords. At a chosen life
fraction we take each engine's \emph{true} $10$-dimensional health vector $\mathbf{x}$
(\cref{ch:trad}), find its $k=10$ nearest neighbours in standardized health space --- engines
in a near-identical present condition --- and record the standard deviation of their
\emph{true} remaining lives. Two engines in the same health state that nonetheless fail
hundreds of cycles apart differ for reasons no present measurement contains; that spread is
aleatoric by construction. The floor is the fleet-median of this local spread. Because it
conditions on \emph{true} health, not an estimate, no better sensor or model can lower it ---
it is a property of the degradation process, not of any method.

\begin{figure}[htbp]
  \centering
  \includegraphics[width=0.72\textwidth]{prognostic_floor}
  \caption{The prognostic floor (tag \texttt{prereg-v11}). Red: the ground-truth aleatoric
  floor --- the RUL spread that survives conditioning on the engine's exact current health.
  Blue: the epistemic gap between the best achieved method (label) and that floor. Early in
  life the floor dominates; late in life a $4\times$ reducible gap opens.}
  \label{fig:prognostic-floor}
\end{figure}

The result (\cref{fig:prognostic-floor}) is a two-part verdict, frozen as
hypotheses \textbf{H11.1} and \textbf{H11.2} before measurement:

\begin{itemize}
  \item \textbf{H11.1 confirmed --- the floor is real and dominant.} At mid-life the
    aleatoric floor is $1053$ cycles against a marginal true-RUL spread of $1205$: fully
    $87\,\%$ of the raw uncertainty is irreducible even with perfect knowledge of the current
    state. Remaining life is, to first order, \emph{not} written in the present --- a sober
    correction to the implicit promise of much prognostics work.
  \item \textbf{H11.2 confirmed --- the reducible headroom concentrates late.} The best
    method's RMSE sits at $1.6\times$ the floor early ($1694$ vs $1053$) but $4.0\times$ late
    ($858$ vs $212$). Early, the method is already near the physical wall and further effort
    buys little; late, most of the error is epistemic and a better prognostic can still cut it
    substantially. The certificate says \emph{where} prognostic investment pays off.
\end{itemize}

Together the two certificates bound both halves of health monitoring against ground truth:
the identifiability certificate says what the present data lets you \emph{diagnose}, and the
prognostic floor says what the present state lets you \emph{predict}. Each replaces a vague
``it's hard'' with a measured, validated number and --- crucially --- with the direction that
number points: which sensor to buy for diagnosis, and at which life stage to invest for
prognosis. The honest boundary applies here too: the floor is estimated by a
finite-neighbourhood oracle, so it is itself an estimate with sampling noise, and it bounds
the \emph{achievable} RMSE, not the error of any one prediction. What it delivers is the
prognostic counterpart of the first certificate's promise --- \emph{here is precisely how much
of this engine's future its present can and cannot reveal}.
